\documentclass[svgnames]{article}     % use "amsart" instead of "article" for AMSLaTeX format
%\geometry{landscape}                 % Activate for rotated page geometry

%\usepackage[parfill]{parskip}        % Activate to begin paragraphs with an empty line rather than an indent

\usepackage{graphicx}                 % Use pdf, png, jpg, or eps§ with pdflatex; use eps in DVI mode

%maths                                % TeX will automatically convert eps --> pdf in pdflatex
\usepackage{amssymb}
\usepackage{amsmath}
\usepackage{esint}
\usepackage{geometry}

% Inverting Color of PDF
%\usepackage{xcolor}
%\pagecolor[rgb]{0.19,0.19,0.19}
%\color[rgb]{0.77,0.77,0.77}

%noindent
\setlength\parindent{0pt}

%pgfplots
\usepackage{pgfplots}

%images
%\graphicspath{{ }}                   % Activate to set a image directory

%tikz
\usepackage{pgfplots}
\pgfplotsset{compat=1.15}
\usepackage{comment}
\usetikzlibrary{arrows}
\usepackage[most]{tcolorbox}

%Figures
\usepackage{float}
\usepackage{caption}
\usepackage{lipsum}

\tcbset {
  base/.style={
    arc=0mm,
    bottomtitle=0.5mm,
    boxrule=0mm,
    colbacktitle=black!10!white,
    coltitle=black,
    fonttitle=\bfseries,
    left=2.5mm,
    leftrule=1mm,
    right=3.5mm,
    title={#1},
    toptitle=0.75mm,
  }
}

\definecolor{brandblue}{rgb}{0.34, 0.7, 1}
\newtcolorbox{mainbox}[1]{
  colframe=brandblue,
  base={#1}
}

\newtcolorbox{subbox}[1]{
  colframe=black!30!white,
  base={#1}
}


\title{The Minimal Polynomial -- Eigenvalue Relation}
\author{Deval Deliwala}
%\date{}                              % Activate to display a given date or no date

\begin{document}
\maketitle
%\section{}
%\subsection{}
%\tableofcontents                     % Activate to display a table of contents


The goal of this note is to derive \textit{what eigenvalues are} without using
the normal method of determinants and characteristic equations. 

\section{The Minimal Polynomial}

\begin{mainbox}{Monic Polynomials}
  A \textit{monic polynomial} is a polynomial whose highest-degree coefficient
  equals 1. 

  \[
  5 + 5x + x^2 
  \] \vspace{3px}
   is a monic polynomial of degree 2. 
  
\end{mainbox}

We begin with the following lemma, 

\begin{subbox}{Existence and Uniqueness of Minimal Polynomials} 
Suppose $V$ is a finite-dimensional vector space and $T \in \mathcal{L}(V)$.
Then, there exists a \textit{unique} monic polynomial $p \in
\mathcal{P}(\mathbb{F})$ of smallest degree such that $p(T) = 0$. Furthermore,
$\text{deg } p \leq \text{dim } V$
\end{subbox} \mbox{} \\

\textbf{\textit{Proof}} \\ 

If  dim $V = 0$, then $I$ is the zero operator on $V$ and thus we take $p$ to
be the constant polynomial 1. \\

Now let's assume dim $V > 0$ and the desired result is true for all operators
on all complex vectors of smaller dimension. Let $v \in V$ be such that $v \neq
0$. The list $v, Tv, \cdots, T^{\text{dim} V} v$ has a length $1 + \text{dim } V$ and
is therefore linearly dependent. \\

By the linear independence lemma, there exists a smallest positive integer $m
\leq \text{dim } V$ such that $T^m v$ is a linear combination of $v, Tv,
\cdots, T^{m-1} v$ before it. Thus there exists scalars $c_0, c_1, c_2, \cdots
, c_{m-1} \in \mathbb{F}$ such that 
\begin{align*}
  c_0v + c_1 Tv + \cdots + c_{m-1}T^{m-1}v &= - T^m v\\
  c_0v + c_1 Tv + \cdots c_{m-1} T^{m-1} v + T^m v &= 0
\end{align*}

We define a monic polynomial $q \in \mathcal{P}_m(\mathbb{F})$ by 

\[
  q(z) = c_0 + c_1z + \cdots + c_{m-1} z^{m-1} + z^m
\] \vspace{3px}

Therefore, substituting in $q(T)$, we see that $q(T) = 0$. 

If a variable  $k$ is some nonnegative integer, then 

\[
q(T) (T^k v) = T^k (q(T) v) = T^k (0) = 0. 
\] \vspace{3px}

And since $v, Tv, \cdots, T^{m-1} v$ is linearly independent, the above
equation implies $\text{dim null} q(T) \geq m$. Hence, 

\[
\text{dim range} q(T) = \text{dim} V - \text{dim null}q(T) \leq \text{dim}
V - m. 
\] \vspace{3px}

And since the range $q(T)$ is invariant under $T$, we apply induction onto the
operator $T\big|_{\text{range} q(T) } $ onto the vector space range $q(T)$. We
fine that there is a monic polynomial $s \in \mathcal{P}(\mathbb{F})$ with 

\[
\text{deg }  s\leq \text{dim } V - m \quad \text{ and } \quad
s(T\big|_{\text{range} q(T)} ) = 0. 
\] \vspace{3px}


Therefore for all $v \in V$, we have 

\[
  (sq)(T)(v) = s(T)(q(T)v) = 0
\] \vspace{3px}

which implies $sq$ is a monic polynomial such that deg $sq \leq$ dim $V$ and
$(sq)(T) = 0$. This polynomial is called the \textit{minimal polynomial}. 

\begin{mainbox}{Minimal Polynomial}
  Suppose  $V$ is finite-dimensional and $T \in \mathcal{L}(V)$. Then the
  \textit{minimal polynomial} of $T$ is the unique monic polynomial $p \in
  \mathcal{P}{\mathbb{F}}$ of smallest degree such that $p(T) = 0$ (the
  0 operator). 
\end{mainbox}

\textit{Example}: \mbox{} \\ 

Suppose $T \in \mathcal{L}(\mathbb{F}^5)$ and 

\[
\mathcal{M} (T) = \begin{pmatrix}
  0 & 0 & 0 & 0 & -3 \\ 
  1 & 0 & 0 & 0 & 6 \\
  0 & 1 & 0 & 0 & 0 \\
  0 & 0 & 1 & 0 & 0 \\
  0 & 0 & 0 & 1 & 0 
\end{pmatrix} 
\] \vspace{3px}

with respect to the standard bases. Taking $v = e_1$, we have 

\begin{align*}
  T e_1 &= e_2, \\
  T^2 e_1 &= T(Te_1) = Te_2 = e_3 \\
  T^3e_1 &= T(T^2 e_1) = Te_3 = e_4 \\ 
  T^4e_1 &= T(T^3 e_1) = Te_4 = e_5 \\ 
  T^{5}e_1 &= T(T^{4}e_1) = Te_5 = -3e_1 + 6e_2 
\end{align*}

Therefore, $3e_1 - 6Te_1 = -T^5 e_1$. Because $e_1, e_2, \hdots, e_5$ is
linearly independent, so too is $e_1, Te_1, \hdots, T^4 e_1$, as they are
equivalent. Thus, no other linear combination of the list equals $-T^5 e_1$.
Hence the minimal polynomial of $T$ is $p(z) = 3 - 6z + z^5$. \\

Now comes the main point of all this. 

\begin{mainbox}{Eigenvalues are the zeros of the minimal polynomial}
  Suppose $V$ is finite dimensional and $T \in \mathcal{L}(V)$.  

  \begin{itemize}
    \item[(a)] The zeros of the minimal polynomial of $T$ are the eigenvalues
      of $T$. 
    \item[(b)] If $V$ is a complex vector space, then the minimal polynomial of
      $T$ has the form 
      \[ (z - \lambda_1) \cdots (z - \lambda_m), \] 
      
      where $\lambda_1, \hdots, \lambda_m$ are the eigenvalues of $T$, possibly
      with repetitions. 
  \end{itemize}
\end{mainbox} \mbox{} \\

\textbf{\textit{Proof}}  \\

(a) \\ 
Let $p$ be a minimal polynomial of $T$. \\

First suppose $\lambda \in \mathbb{F}  $ is a zero of $p$. Then $p$ can be
written in the form 

\[
p(z) = (z - \lambda)q(z), 
\] \vspace{3px}

where $q(z)$ is a monic polynomial with coefficients in $\mathbb{F}  $. Because
$p(T) = 0$, we get 

\[
0 = (T - \lambda I )(q(T) v) 
\] \vspace{3px}

for all $v\in V$. And since the degree of $q$ is less than the degree of the
minimal polynomial $p$, there is at least one vector $v \in V$ such that
$q(T)v \neq 0$. The equation above the implies $\lambda$ is an eigenvalue of
$T$. 

To prove that every eigenvalue of $T$ is a zero of $p$, now suppose $\lambda
\in \mathbb{F}  $ is an eigenvalue of $T$. Then there exists $v\in V$ such that
$Tv  = \lambda v$. Repeated applications of $T$ to both sides show that $T^k
v = \lambda^k v$ for every nonnegative $k$. Therefore, 

\[
p(T) v = p(\lambda) v
\] \vspace{3px}

And since $p$ is the minimal polynomial, we have $p(T) v = 0 = p(\lambda)v$,
implying $p(\lambda) = 0$. Therefore $\lambda$ is a zero of $p$. \\

We can prove $(b)$ using the above result and the fundamental theorem of
algebra. 

The next result completely characterizes the polynomials that, when applied to
an operator, give the 0 operator.  

\begin{mainbox}{$q(T) = 0\Leftrightarrow q$ is a polynomial multiple of the
  minimal polynomial}
  Suppose $ V$ finite dimensional, $T \in \mathcal{L} (V)$ and $q \in
  \mathcal{P} (\mathbb{F} ) $. Then $q(T) = 0$ if and only if $q$ is
  a polynomial multiple of the minimal polynomial of $T$. 
\end{mainbox} \mbox{} \\

\textbf{\textit{Proof}}\\

\textbf{Forward Direction} (Assume $q(T) = 0$, prove $q$ multiple of $p$ ). \\
The
polynomial division algorithm states there exist polynomials $s, r \in
\mathcal{P} (\mathbb{F} ) $ such that 

\[
q = ps + r
\] \vspace{3px}

and deg $r$ $<$ deg $p$. Applying $T$ to the equation, we get 

\[
q(T) = p(T) s(T) + r(T)
\] \vspace{3px}

But since $p$ is the minimal polynomial, we know $p(T) = 0 $. Therefore, this
reduces to 

\[
q(T) = r(T) = 0
\] \vspace{3px}

which implies $r(T) = 0$. And since the minimal polynomial $p$ is
\textit{defined} as the polynomial of smallest degree such that $p(T) = 0$, $r$
therefore must equal 0. Hence $q = ps$: a polynomial multiple of $p$. \\


\textbf{Backward Direction} (Assume $q$ a polynomial multiple of $p$, prove
$q(T) = 0$ ). \\ 

This direction is must easier, we just apply $T$ to the equation 

\begin{align*}
  q &= ps  \\ q(T) &= p(T)s(T) \\ q(T) &= 0s(T) \\ q(T) &= 0 
\end{align*}

An even cooler result that  follows is 

\begin{mainbox}{Minimal polynomial of a restriction operator}
  
  Suppose $V$ is finite-dimensional, $T \in \mathcal{L} (V)$ and $U$ is
  a subspace of $V$ that is invariant under $T$. Then the minimal polynomial of
  $T$ is a polynomial multiple of the minimal polynomial of $T\big|_U$. 

\end{mainbox}

The next result is also important, showing that the constant term of the
minimal polynomial of an operator can show whether the operator is invertible! 


\begin{mainbox}{$T$ not invertible $\Leftrightarrow$ constant term of minimal
polynomial of $T$ is 0}
  Suppose $V$ is finite-dimensional and $T\in \mathcal{L} (V)$. Then $T$ is not
  invertible if and only if the constant term of the minimal polynomial of $T$ 
  \textit{is} 0. 
\end{mainbox}


The proof is fairly simple: 

\begin{align*}
  T \text{ is not invertible}  &\Longleftrightarrow 0 \text{ is an eigenvalue
  of }  T \\ &\Longleftrightarrow 0 \text{ is a zero of } p \\ &\Longleftrightarrow \text{the constant term of } p \text{ is 0} 
\end{align*}
\end{document}
